\section{Results}
\label{sec:results}

We analyzed data collected from 4,300 repositories, with a total of 4,321 tool-related observations, including 11,098 rule files and 86,339 recorded edits.
Research findings are organized by research question. For each question, we first present the main results and figures. Next, we examine the differences between tools. Finally, we assess whether the findings support the corresponding research hypothesis.

\subsection{RQ1: To what extent do repositories expose evidence of architecture-checker invocation in GitHub Actions workflows?}
\label{sec:rq1}

To answer RQ1, we examined the default-branch GitHub Actions
workflows of repositories using import-linter and
dependency-cruiser. We hypothesized that fewer than half of the
repositories using each tool would invoke it in CI.
Table~\ref{tab:ci_results} show
the proportions of repository-tool observations meeting this
criterion.

For import-linter, we found that 351 out of 826 observations satisfy
the integration criterion (42.5\%, 95\% CI: 39.4-45.8\%). For
dependency-cruiser, we found that only 125 out of 1,073 observations
meet the criterion (11.6\%, 95\% CI: 9.8-13.5\%). Both proportions
are below 50\%, which supports our hypothesis for these two tools.
We also combined the observations from both tools and found that
476 out of 1,899 observations satisfy the criterion (25.1\%, 95\%
CI: 23.1-27.1\%). This result is also below the hypothesized
threshold of 50\%.

In ArchUnit observations, 869 out of 2,421 with an available workflow result contain at least one
of these commands (35.9\%, 95\% CI: 34.1-37.9\%). Because ArchUnit
does not provide a dedicated CI command, this result only shows that
the repository invokes a potentially relevant Maven or Gradle build
or test task. It does not confirm that ArchUnit rules are actually
executed. Therefore, we do not use this result to evaluate the hypothesis.

The observed workflow-invocation rates differ across the two tools
with dedicated commands: import-linter's rate is approximately
3.7 times that of dependency-cruiser. This contrast should be
interpreted in light of the tool-specific detection criteria rather
than as a direct comparison of equivalent enforcement mechanisms.
In particular, the ArchUnit estimate is a build-or-test proxy and
is not directly comparable to the dedicated-command measures.

\begin{table}[t]
\centering
\caption{Architecture-checker invocation evidence in GitHub Actions by tool.}
\label{tab:ci_results}
\setlength{\tabcolsep}{10pt}
\setlength{\aboverulesep}{0pt}
\setlength{\belowrulesep}{0pt}
\renewcommand{\arraystretch}{1.25}
\begin{tabular}{@{}lrrrr@{}}
\toprule
\textbf{Tool} & \textbf{Criterion met} & \textbf{$N$} &
\textbf{Share} & \textbf{95\% CI} \\
\midrule
\rowcolor{rowband}
ArchUnit            & 869 & 2{,}421 & 35.9\% & 34.1-37.9 \\
import-linter       & 351 & 826     & 42.5\% & 39.4-45.8 \\
\rowcolor{rowband}
dependency-cruiser  & 125 & 1{,}073 & 11.6\% &  9.8-13.5 \\
\bottomrule
\end{tabular}
\end{table}

%\begin{figure}[t]
\centering
\begin{tikzpicture}
\begin{axis}[
  xbar,
  width=0.95\columnwidth,
  height=4.2cm,
  xmin=0, xmax=65,
  xlabel={Share in GitHub Actions (\%)},
  symbolic y coords={dependency-cruiser,ArchUnit proxy,import-linter},
  ytick=data,
  nodes near coords,
  nodes near coords align={horizontal},
  every node near coord/.append style={font=\scriptsize},
  bar width=10pt,
  enlarge y limits=0.28,
  grid=major,
  major grid style={gray!25},
  tick label style={font=\scriptsize},
  label style={font=\scriptsize},
]
\addplot[fill=teal!70!black, draw=teal!80!black] coordinates {
  (42.5,import-linter)
  (35.9,{ArchUnit proxy})
  (11.6,dependency-cruiser)
};
\draw[dashed, red!70!black, thick]
  (axis cs:50,dependency-cruiser) -- (axis cs:50,import-linter)
  node[above, font=\footnotesize, red!70!black, midway, yshift=3pt] {$50\%$};
\end{axis}
\end{tikzpicture}
\caption{GitHub Actions integration rates by tool. The ArchUnit bar represents a Maven or Gradle build-or-test proxy rather than direct evidence of ArchUnit rule execution. The dashed line indicates the hypothesized threshold of 50\%.}
\label{fig:ci_share}
\end{figure}

\begin{keyfinding}
Our results show that fewer than half of the repositories using
import-linter and dependency-cruiser invoke these tools in GitHub
Actions workflows. This supports our hypothesis for both tools.
For ArchUnit, we report it separately because the Maven or Gradle
build-or-test commands we identified do not provide direct evidence
that ArchUnit rules are executed.
\end{keyfinding}


\subsection{RQ2: Which architecture constraint dominates executable rules, and how do they differ across tools?}
\label{sec:rq2}
To answer RQ2, we classified 21,713 held-out constructs using our
predefined taxonomy, which applies the same categories to all three
tools (Table~\ref{tab:rule_taxonomy}). Our parser successfully
categorized 15,173 of 17,907 ArchUnit constructs (84.7\%),
2,136 of 2,265 import-linter constructs (94.3\%), and
1,503 of 1,541 dependency-cruiser constructs (97.5\%). The remaining
2,901 constructs were not categorized and were excluded from the
category proportions. Figure~\ref{fig:categories_by_tool} shows
the distribution of the categorized constructs for each tool.


We found that dependency and layering constraints account for
98.1\% of categorized import-linter constructs (95\% CI:
97.5-98.6\%). For ArchUnit and dependency-cruiser, these
constraints account for 58.5\% (95\% CI: 57.8-59.3\%) and
56.8\% (95\% CI: 54.4-59.1\%), respectively. Only the
import-linter confidence interval lies entirely above 70\%.
Therefore, our results support the hypothesis for import-linter,
but not for ArchUnit or dependency-cruiser.

We also found differences in the other constraint categories.
For dependency-cruiser, cycle constraints account for 649 of
1,503 categorized constructs (43.2\%). For ArchUnit, naming and
annotation constraints account for 18.5\% and 8.4\%,
respectively. Visibility and other constraints together account
for an additional 11.5\%. We further compared the overall category
distributions and found a difference across the three tools
($\chi^2 = 5865.56$, $df = 12$, Cram\'er's $V = 0.395$).

These results show that the three tools have different constraint
profiles. Import-linter focuses almost entirely on dependency and
layering constraints. Dependency-cruiser also covers these
constraints, but a substantial proportion of its constructs focus
on cycle constraints. In contrast, ArchUnit covers a wider range
of constraint types, including naming and annotation rules.
Therefore, the lower proportion of dependency and layering
constraints in ArchUnit does not necessarily mean that these
constraints are less important in its rules. Instead, it reflects
a more diverse distribution of constraint categories.

\begin{figure}[t]
\centering
\fbox{\parbox{0.9\columnwidth}{
\textbf{Figure omitted in source manuscript.}\newline
Placeholder for category-level distribution of rule semantics by tool.
}}
\caption{Distribution of executable rule categories across ArchUnit, import-linter, and dependency-cruiser.}
\label{fig:categories_by_tool}
\end{figure}


\begin{keyfinding}
Our results show that dependency and layering constraints
dominate import-linter (98.1\%), but account for smaller shares
of ArchUnit (58.5\%) and dependency-cruiser (56.8\%). Therefore,
our hypothesis is supported only for import-linter. We also
found that the distribution of constraint categories differs
across the three tools (Cram\'er's $V = 0.395$).
\end{keyfinding}


\subsection{RQ3: How often are architecture rules exempted, disabled, or left unchanged?}
\label{sec:rq3}

To answer RQ3, we analyze the historical edits of
11,098 rule files and classify each edit according
to its effect on rule enforcement. We then identify
the first detected exemption or disabling event
in each file's history and estimate its 24-month
cumulative incidence using the Kaplan-Meier estimator.

Across all tools, we found that the estimated 24-month cumulative incidence
of rule rot events is 26.3\% (95\% CI: 25.4-27.2\%). This estimate
is below our hypothesized threshold of 30\%. However, the
results differ considerably across tools. ArchUnit has the
lowest estimated incidence at 12.9\% (8,579 files; 95\% CI:
12.1-13.8\%). In contrast, the estimated incidence is 58.1\%
for import-linter (1,093 files; 95\% CI: 54.0-62.3\%) and
80.5\% for dependency-cruiser (1,426 files; 95\% CI:
78.4-82.7\%). These results show that the overall estimate
hides substantial differences between the three tools.
While ArchUnit remains below the 30\% threshold, both
import-linter and dependency-cruiser exceed it.

We also examine the proportion of rule files with no
subsequent modifications. We find that 5,501 of the 
11,098 files (49.6\%, 95\% CI: 48.6-50.4\%) have only 
one observed commit in their histories. We calculate
this confidence interval using percentile bootstrap 
resampling at the file level with 1,000 replicates. 
This measure differs from rule-rot incidence because 
it considers all subsequent modifications rather than 
only edits classified as exemption or disabling events. 
A file can therefore receive subsequent modifications 
without experiencing a candidate rule-rot event.

Among 5,122 files with at least 24 months of follow-up,
49.3\% receive no subsequent modifications (95\% CI:
48.0-50.7\%). We report this result separately because it
is based only on files with sufficient follow-up time. In
contrast, our primary Kaplan-Meier analysis accounts for
incomplete follow-up through right-censoring. These estimates
therefore describe related but different aspects of rule-file
maintenance.

We further examine the types of edits recorded in our dataset.
Table~\ref{tab:edit_class_counts} presents the distribution
of the automatically assigned edit classes. The
\emph{restructure} class is the most common, accounting for
69,047 of the 86,339 edits (80.0\%). Exemption and disabling
edits together account for 4.7\% of all recorded edits.
Although these edits represent a relatively small proportion
of all changes, the first qualifying exemption or disabling
edit in each file defines the event in our primary
Kaplan-Meier analysis.

As a supplementary analysis, we examine an alternative event
definition based on the first \emph{loosen} or \emph{exempt}
edit, excluding \emph{disable}. Using a separate Kaplan-Meier
analysis with this event definition, we estimate a 24-month
cumulative incidence of 41.9\% (95\% CI: 40.9-43.0\%).
This estimate reflects a different aspect of rule-file
evolution and is not directly comparable to the primary
candidate rule-rot incidence.

We also fit a Cox proportional-hazards model using the same
alternative event definition and time-to-event variable.
The model reports a hazard ratio of 1.307 (95\% CI:
1.270-1.345) for log file-level commit count and 0.850
(95\% CI: 0.793-0.911) for log author count. We interpret
these estimates as supplementary associations with the
hazard of the alternative event, rather than as adjusted
estimates of the primary candidate rule-rot outcome.
Because both covariates are calculated over each file's
entire observed history, they may include activity occurring
after the event. We therefore do not interpret these
associations as evidence that file activity or author count
predicts the occurrence of the event.

Overall, our results do not support the hypothesis that the
24-month cumulative incidence of candidate rule-rot events
exceeds the 30\% threshold across the entire dataset.
However, the tool-specific estimates reveal substantial
differences across ecosystems. We also find that nearly
half of the rule files receive no subsequent modifications
during their observed histories. Nevertheless, the absence
of subsequent modifications does not necessarily indicate
that a file has never experienced rule rot, just as the
presence of modifications does not necessarily indicate
rule rot. Most recorded edits are classified as
\emph{restructure}, and our automatic edit classification
has not been validated against manually annotated examples.
We therefore interpret these results in light of the
limitations of our event-detection procedure.

\begin{table}[t]
\centering
\caption{Rule rot and never-modified shares by tool. Rule rot is the Kaplan--Meier incidence of a first exemption or disabling within 24 months.}
\label{tab:rule_rot_by_tool}
\small
\setlength{\tabcolsep}{7pt}
\renewcommand{\arraystretch}{1.25}
\begin{tabular}{lrrr}
\toprule
\textbf{Tool} & \textbf{Files} & \textbf{24-mo rot} & \textbf{Never mod.} \\
\midrule
\rowcolor{rowband}
ArchUnit & 8{,}579 & 12.9\% & 55.7\% \\
import-linter & 1{,}093 & 58.1\% & 13.9\% \\
\rowcolor{rowband}
dependency-cruiser & 1{,}426 & 80.5\% & 40.2\% \\
\midrule
All tools & 11{,}098 & 26.3\% & 49.6\% \\
\bottomrule
\end{tabular}
\end{table}
\begin{figure}[t]
\centering
\begin{tikzpicture}
\begin{axis}[
  ybar,
  width=\columnwidth,
  height=6.4cm,
  ymin=0, ymax=100,
  ylabel={24-month rule-rot incidence (\%)},
  symbolic x coords={ArchUnit,import-linter,dependency-cruiser},
  xtick=data,
  xticklabels={ArchUnit,{import-\\linter},{dependency-\\cruiser}},
  xticklabel style={font=\footnotesize, align=center},
  nodes near coords,
  every node near coord/.append style={font=\footnotesize},
  bar width=18pt,
  enlarge x limits=0.28,
  tick label style={font=\footnotesize},
  label style={font=\footnotesize},
  grid=major,
  major grid style={gray!20},
  error bars/y dir=both,
  error bars/y explicit,
]
\addplot[fill=blue!70!black, draw=blue!80!black, error bars/.cd, y dir=both, y explicit]
  coordinates {
    (ArchUnit,12.9) +- (0,0.85)
    (import-linter,58.1) +- (0,4.15)
    (dependency-cruiser,80.5) +- (0,2.15)
  };
\draw[dashed, red!70!black, thick]
  (current axis.west |- {axis cs:ArchUnit,30}) --
  (current axis.east |- {axis cs:ArchUnit,30})
  node[above, font=\footnotesize, red!70!black, pos=0, anchor=south west, xshift=2pt, yshift=1pt] {$30\%$};
\end{axis}
\end{tikzpicture}
\caption{Kaplan-Meier survival analysis of rule rot by tool.}
\label{fig:survival_analysis}
\end{figure}

\begin{table}[t]
\centering
\caption{Automatic edit-class counts over 86{,}339 revisions of 11{,}098 rule files.}
\label{tab:edit_class_counts}
\small
\setlength{\tabcolsep}{4pt}
\begin{tabular}{@{}lrr@{}}
\toprule
\textbf{Class} & \textbf{Edits} & \textbf{Share} \\
\midrule
Restructure & 69{,}047 & 80.0\% \\
Add & 8{,}448 & 9.8\% \\
Loosen & 3{,}534 & 4.1\% \\
Exempt & 3{,}219 & 3.7\% \\
Tighten & 1{,}187 & 1.4\% \\
Disable & 904 & 1.0\% \\
\bottomrule
\end{tabular}
\end{table}


\begin{keyfinding}
Our results show that the pooled 24-month incidence of
rule rot is 26.3\%, below the hypothesized threshold of
30\%. However, this overall estimate hides substantial
differences across tools, ranging from 12.9\% for ArchUnit
to 80.5\% for dependency-cruiser. We also found that
49.6\% of rule files receive no subsequent modifications.
This measure is distinct from rule-rot incidence because
not all modifications involve exempting or disabling rules.
\end{keyfinding}


\subsection{RQ4: How do dependency and layering violations change around the adoption of ArchUnit rules?}
\label{sec:rq4}

To answer RQ4, we examined how dependency and layering violations
change around the adoption of ArchUnit rules. We defined adoption
as the first commit that introduces a parseable ArchUnit rule.
Our initial comparison set contains 2,422 tool-specific reepository
observations. Among these, we identified 317 repositories with
at least 12 months of observable history both before and after
the adoption commit. Then, we use these 317 repositories to construct the 
longitudinal analysis cohort (Table~\ref{tab:rq4_cohort}). 
First, we exclude 169 repositories without dependency or 
layering constraints that can be evaluated through our 
replay procedure, leaving 148 repositories. We further 
exclude seven repositories with incomplete violation 
counts at one or more of the three observation points 
($t={-}1$, $t=0$, and $t=12$). This results in a final 
cohort of 141 repositories with at least one evaluable 
constraint and complete observations at all three time 
points. This selection ensures that an empty set of 
evaluable constraints is not interpreted as evidence 
of zero architectural violations.

\begin{table*}[t]
\centering
\caption{Analysis sets within the ArchUnit trajectory cohort.
The ratifier set and the evaluable-rule set are separate
filters of the same 317 repositories. The complete-trajectory
set is the subset of the evaluable-rule set with all three
snapshots.}
\label{tab:rq4_cohort}
\small
\setlength{\tabcolsep}{8pt}
\renewcommand{\arraystretch}{1.15}
\begin{tabular}{@{}lrrl@{}}
\toprule
\textbf{Set} & \textbf{$N$} & \textbf{Excluded} & \textbf{Exclusion reason} \\
\midrule
Trajectory cohort & 317 & 2{,}105 of 2{,}422 &
Fewer than 12 months of history before or after adoption \\
Ratifier/reformer & 298 & 19 of 317 &
No violation count at $t{=}{-}1$ or $t{=}0$ \\
Evaluable rules & 148 & 169 of 317 &
No translated dependency or layering rule \\
Complete trajectory & 141 & 7 of 148 &
Missing a count at $t{=}{-}1$, $t{=}0$, or $t{=}12$ \\
\bottomrule
\end{tabular}
\end{table*}


We focus on dependency and layering constraints that 
can be evaluated using reconstructed source-level import 
graphs. The replay procedure does not execute ArchUnit 
itself and therefore covers only the subset of constraints 
that can be translated into checks over these graphs.
We examine the violation trajectories of the 141 
repositories with complete observations. 
Table~\ref{tab:rq4_trajectories} and 
Figure~\ref{fig:rq4_change} present the results. 
The median violation count is zero at all three 
observation points: one month before the introduction 
commit, at the introduction commit, and 12 months 
afterward. The interquartile range is 0-11 at each point. 
In addition, 83 of the 141 repositories already have zero 
evaluable violations one month before the introduction 
commit.

The mean violation counts are 73.3 (95\% CI: 23.9-138.2) 
before introduction, 85.6 (95\% CI: 28.6-154.8) at 
introduction, and 103.9 (95\% CI: 33.8-196.2) 12 months 
afterward. Comparing the first and last observation points, we 
find that violations decrease in 18 repositories 
(12.8\%, 95\% CI: 7.1-17.7\%), remain unchanged in 86 repositories 
(61.0\%, 95\% CI: 53.2-69.5\%), and increase in 37 repositories 
(26.2\%, 95\% CI: 19.1-33.3\%).

We also use within-repository ordinary least squares (OLS) 
regression to estimate violation trends before and after 
the introduction commit. The mean pre-introduction slope 
is 1.63 violations per month (95\% CI: -0.09 to 4.38), 
while the mean post-introduction slope is 1.83 violations 
per month (95\% CI: 0.47 to 3.74). These results do not 
indicate a consistent decline in violations following the 
introduction of ArchUnit rules. Instead, the mean 
violation counts increase across the three observation
points, and the estimated post-introduction slope is 
positive. The median remains zero and the interquartile 
range remains 0-11, indicating a right-skewed distribution 
in which relatively large violation counts in some 
repositories contribute substantially to the mean.

Overall, most repositories show no change in violation 
counts between one month before the introduction commit 
and 12 months afterward. More repositories experience an 
increase than a decrease in violations. Furthermore, 83 
repositories already have zero evaluable violations before 
the introduction of ArchUnit rules. These findings 
describe violation trajectories in repositories that 
introduced parseable ArchUnit constraints and met our 
history and evaluability requirements. However, we cannot 
attribute these changes to ArchUnit adoption because our 
analysis does not include an untreated comparison group. 
The results should therefore be interpreted as descriptive 
patterns rather than causal effects.

\begin{table}[t]
\centering
\caption{Dependency and layering violation counts at $t={-}1$, $t=0$, and $t=12$ months relative to the introduction of parseable ArchUnit constraints.}
\label{tab:rq4_trajectories}
\small
\setlength{\tabcolsep}{3.5pt}
\setlength{\aboverulesep}{0pt}
\setlength{\belowrulesep}{0pt}
\renewcommand{\arraystretch}{1.25}
\begin{tabular}{@{}lrrrr@{}}
\toprule
\textbf{Point} & \textbf{Median} & \textbf{IQR} & \textbf{Mean} & \textbf{95\% CI (mean)} \\
\midrule
\rowcolor{rowband}
$t{=}{-}1$ & 0 & 0--11 & 73.3 & 23.9--138.2 \\
$t{=}0$ & 0 & 0--11 & 85.6 & 28.6--154.8 \\
\rowcolor{rowband}
$t{=}12$ & 0 & 0--11 & 103.9 & 33.8--196.2 \\
\bottomrule
\end{tabular}
\end{table}

\begin{figure}[t]
\centering
\begin{tikzpicture}
\begin{axis}[
  ybar,
  width=\columnwidth,
  height=5.6cm,
  ymin=0, ymax=78,
  ylabel={Share of series (\%)},
  symbolic x coords={Decreased,Unchanged,Increased},
  xtick=data,
  nodes near coords,
  every node near coord/.append style={font=\footnotesize},
  bar width=24pt,
  enlarge x limits=0.25,
  tick label style={font=\footnotesize},
  label style={font=\footnotesize},
  grid=major,
  major grid style={gray!20},
]
\addplot[fill=teal!70!black, draw=teal!80!black] coordinates {
  (Decreased,12.8)
  (Unchanged,61.0)
  (Increased,26.2)
};
\end{axis}
\end{tikzpicture}
\caption{Change in dependency and layering violations over 12 months after the introduction of parseable ArchUnit constraints.}
\label{fig:rq4_change}
\end{figure}


\begin{keyfinding}
Among the 141 repositories with complete violation trajectories, 61.0\% show no change in violation counts between one month before the introduction of ArchUnit rules and 12 months afterward. More repositories experience an increase (26.2\%) than a decrease (12.8\%). The median violation count remains zero across all three observation points, while the mean increases from 73.3 before introduction to 103.9 after 12 months. These findings describe observed violation trajectories but do not establish a causal relationship between ArchUnit adoption and changes in architectural violations.
\end{keyfinding}